% Generated draft: requires mathematical and source review.
\documentclass[11pt]{article}
\usepackage[margin=25mm]{geometry}
\usepackage{fontspec}
\setmainfont{Latin Modern Roman}
\usepackage{amsmath,amssymb,mathtools,unicode-math}
\setmathfont{Latin Modern Math}
\usepackage{enumitem,needspace,xurl,hyperref}
\setlength{\parindent}{0pt}
\setlength{\parskip}{5pt}
\setlength{\emergencystretch}{3em}
\newcommand{\N}{\mathbb{N}}
\newcommand{\Z}{\mathbb{Z}}
\newcommand{\Q}{\mathbb{Q}}
\newcommand{\R}{\mathbb{R}}
\newcommand{\C}{\mathbb{C}}
\newcommand{\F}{\mathbb{F}}
\newcommand{\A}{\mathbb{A}}
\newcommand{\PP}{\mathbb{P}}
\newcommand{\E}{\mathbb{E}}
\newcommand{\eps}{\varepsilon}
\newcommand{\dd}{\,\mathrm{d}}
\begin{document}
\section*{UnsolvedMath ID: 800285}
\textbf{Source problem code:} catalogue entry
\subsection{Definitions and mathematical statement}
Let \(\Delta=\{\zeta\in\C:|\zeta|<1\}\), with Poincar\'e distance
\[
\rho_\Delta(a,b)=\operatorname{arctanh}\left|\frac{a-b}{1-\overline a b}\right|.
\]
In the usual domain interpretation, \(\Omega\subset\C^n\) is a connected open set, where \(n\geq1\). For \(z\in\Omega\) and \(v\in\C^n\), define the infinitesimal Kobayashi pseudometric by
\[
\kappa_\Omega(z;v)=\inf\{ |\alpha|:\alpha\in\C,\ f:\Delta\to\Omega\text{ holomorphic},\ f(0)=z,\ \alpha f'(0)=v\}.
\]
Its integrated pseudodistance is
\[
d_\Omega^{\mathrm{int}}(p,q)=\inf_\gamma\int_0^1\kappa_\Omega\bigl(\gamma(t);\gamma'(t)\bigr)\,\dd t,
\]
where the infimum is over piecewise \(C^1\) paths \(\gamma:[0,1]\to\Omega\) with \(\gamma(0)=p\) and \(\gamma(1)=q\). The source calls this the integrated Kobayashi distance.

A Kobayashi chain of \(m\) discs from \(p\) to \(q\) consists of holomorphic maps \(f_j:\Delta\to\Omega\) and points \(a_j,b_j\in\Delta\), for \(1\leq j\leq m\), such that
\[
f_1(a_1)=p,\qquad f_j(b_j)=f_{j+1}(a_{j+1})\quad(1\leq j<m),\qquad f_m(b_m)=q.
\]
Its length is \(\sum_{j=1}^m\rho_\Delta(a_j,b_j)\).

Under the usual \(C^2\) boundary convention, strong pseudoconvexity means that near every boundary point there is a real \(C^2\) defining function \(r\), with \(\Omega=\{r<0\}\) locally and \(\dd r\ne0\) on the boundary, whose Levi form satisfies
\[
\sum_{j,k=1}^n\frac{\partial^2 r}{\partial z_j\partial\overline z_k}(x)v_j\overline v_k>0
\]
for every nonzero \(v\in\C^n\) satisfying \(\sum_{j=1}^n r_{z_j}(x)v_j=0\), at each boundary point \(x\). This convention explains the terminology; the supplied source does not specify boundary regularity, boundedness, or ambient dimension.

\textbf{Requested assertion, subject to the unresolved scope qualifications below.} Prove that strong pseudoconvexity guarantees an absolute integer \(K\geq1\) such that for every \(p,q\in\Omega\) there is a Kobayashi chain of \(m\leq K\) discs from \(p\) to \(q\) satisfying the exact equality
\[
d_\Omega^{\mathrm{int}}(p,q)=\sum_{j=1}^m\rho_\Delta(a_j,b_j).
\]
Here \(K\) must be independent of \(p,q\). The word ``absolute'' suggests further uniformity, but the source does not resolve whether \(K\) is independent of \(\Omega\), of \(n\), or both. The displayed assertion therefore requires a human determination of its intended domain class and quantifier scope. Exact realization is the literal reading of ``realized''; replacing equality by arbitrarily close approximation would be a different formulation. Using at most \(K\) discs expresses the bounded chain requirement; chains can be padded to exactly \(K\) discs using constant maps. The source also states parenthetically that on a convex domain one can take \(K=1\), without specifying additional hypotheses for that comparison.
\subsection{Short English statement}
The problem asks whether strong pseudoconvexity forces every integrated Kobayashi distance to be attained by a chain with a uniformly bounded number of holomorphic discs. The source leaves the precise domain hypotheses and the scope of the proposed absolute bound unresolved. Its convex-domain comparison states \(K=1\). This is a statement-only draft requiring mathematical review.
\subsection{Sources}
\begin{itemize}
\item Local registry: \texttt{data/unsolvedmath/problems.json} and, for appended records, \texttt{data/unsolvedmath/additions/}; supplied complete record with permanent UnsolvedMath ID \(800285\), exact source problem code \texttt{catalogue entry}, and upstream revision \texttt{d371ee917ac9647fbb06cd6fc529deec6f3dda82}.
\item Supplied reference: Buzzard, Fornaess, and Krantz, \emph{The Several Complex Variables Problem List} (1995), Problem 85, \url{https://arxiv.org/abs/math/9509202}. The registry attributes the problem to Buzzard, Fornaess, and Krantz, with proposed year 1995.
\item Supplied provenance records access on 2026-10-16 and extraction by \texttt{arxiv\_latex}. These are supplied provenance claims; no independent source visit or verification is asserted here.
\item The source offers a prize of \$50, attributed to S. Krantz, and gives the contact text \texttt{sk @math.wustl.edu}. The present availability of the prize and contact has not been verified.
\item The supplied registry labels the problem \texttt{open}. This is a source status label, not an independently verified literature assessment.
\end{itemize}
\end{document}
